On FetchReach (Figure~\ref{fig:frontier}, Table~\ref{tab:rl}), the benefit of ErrP feedback rises
monotonically with decoder quality: across all runs, the paired gain in normalized success AUC over the sparse
run of the same seed is rank-correlated with feedback AUROC (Spearman $\rho=\RLRhoReach{}$, $p=\RLRhoPReach{}$).
Exact feedback raises the AUC from \RLAUCReachSparse{} to \RLAUCReachExact{} (\RLPosReachExact{} seeds better,
sign test $p=\RLSignPReachExact{}$), about half of the gain of the dense reward (\RLAUCReachDense{}). The mean gain
crosses zero at AUROC \RLBreakEvenReach{} (linear interpolation between the measured levels). Below this point the
feedback is harmful: at AUROC 0.55 the AUC drops to \RLAUCReachLow{} (\RLGainReachLow{}, \RLPosReachLow{} seeds
better, $p=\RLSignPReachLow{}$). Final success saturates for every condition (at least \RLSuccReachLow{}), so the
effect is on learning speed. Noisy feedback is therefore not merely uninformative; a likely reason is that, with a
fixed weight, it adds a reward term that is only weakly correlated with progress and that the critic must
fit nonetheless.

On FetchPush, SAC without hindsight relabelling does not learn the task within 500k steps in any condition
(final success \RLSuccPushSparse{} with sparse reward, \RLSuccPushDense{} with dense reward,
\RLSuccPushExact{} with exact ErrP feedback), and no condition differs reliably from sparse. Push therefore shows
that ErrP feedback of any quality does not substitute for exploration when the learner rarely moves the object;
it does not locate a frontier. The staged label rule, which provides several times more feedback events, gives
the same picture (Table~\ref{tab:rl}, right).

Placing the benchmark on the Reach frontier (vertical lines in Figure~\ref{fig:frontier}): the best
cross-subject decoders of cursor errors, cursor-monitoring errors, response errors and orthosis errors lie above
the break-even point, whereas the best decoder of observed robot errors (\BestErrpHriRobot{}) lies well below
it, in the range where feedback slowed learning.
